\section*{الفصل \ref{ch:b2:solid-liquid-diagrams} --- المخططات الثنائية صلب--سائل}

\begin{solution}{exo:b2:solid-liquid-diagrams:1}
عند \qty{1300}{\degreeCelsius} يقع \omterm{def:b2:solid-liquid-diagrams:liquidus}{خط السيولة} عند 0.541 \omterm{def:b2:solid-liquid-diagrams:liquidus}{وخط الصلابة} عند 0.609.
0.40: سائل فقط. 0.58: سائل (0.541) \omterm{def:b2:solid-liquid-diagrams:solid-solution}{ومحلول صلب} (0.609). 0.70:
\omterm{def:b2:solid-liquid-diagrams:solid-solution}{محلول صلب} فقط.
\end{solution}

\begin{solution}{exo:b2:solid-liquid-diagrams:2}
كسر الجسم الصلب $(0.58 - 0.541)/(0.609 - 0.541) = 0.57$: \qty{1.15}{mol} من الجسم الصلب،
و\qty{0.85}{mol} من السائل. النيكل: $1.15 \times 0.609 = \qty{0.70}{mol}$ في
الجسم الصلب، و$0.85 \times 0.541 = \qty{0.46}{mol}$ في السائل؛ والمجموع \qty{1.16}{mol}
$= 2.00 \times 0.58$.
\end{solution}

\begin{solution}{exo:b2:solid-liquid-diagrams:3}
عند ضغط ثابت $v = n + 1 - \varphi$ (لا تفاعل). (أ) $1 + 1 - 2 = 0$.
(ب) $2 + 1 - 2 = 1$. (ج) سائل \omterm{def:b2:solid-liquid-diagrams:solid-solution}{ومحلولان صلبان}: $2 + 1 - 3 = 0$.
\end{solution}

\begin{solution}{exo:b2:solid-liquid-diagrams:4}
انخفاض منتظم حتى \qty{232}{\degreeCelsius}، ثم استواء أفقي بينما يتجمد
القصدير، ثم انخفاض منتظم من جديد، أبطأ قرب درجة حرارة الغرفة. وعلى
الاستواء يتعايش السائل والجسم الصلب عند ضغط ثابت، \omterm{def:b2:equilibrium-shifts:variance}{بتغايرية} 0: فالحرارة
المنزوعة يوفرها التبلور، ولا يمكن أن تتغير درجة الحرارة
حتى يتجمد آخر سائل.
\end{solution}

\begin{solution}{exo:b2:solid-liquid-diagrams:5}
البنزين: $\ln x_B = -(9870/8.314)(1/269.6 - 1/278.64) = -0.1429$، $x_B =
0.867$. النفثالين: $\ln x_N = -(19\,100/8.314)(1/269.6 - 1/353.2) = -2.017$،
$x_N = 0.133$. والمجموع 1.000: فدرجة الحرارة \qty{269.6}{K} ($\qty{-3.5}{\degreeCelsius}$) هي
درجة الحرارة الأوتكتيكية، وفي السائل الأوتكتيكي 0.133 من النفثالين.
\end{solution}

\begin{solution}{exo:b2:solid-liquid-diagrams:6}
يبرد السائل حتى \omterm{def:b2:solid-liquid-diagrams:liquidus}{خط السيولة} من جهة الرصاص، حيث تظهر بلورات (Pb)؛ ويغتني
السائل بالقصدير، نزولًا حتى الأوتكتيك عند \qty{182.2}{\degreeCelsius}
و62.1~\%. وهناك يتجمد السائل المتبقي إلى \omterm{def:b2:solid-liquid-diagrams:eutectic}{الخليط الأوتكتيكي}.
والكسر الكتلي للأوتكتيك: $(30 - 18.91)/(62.13 - 18.91) = 0.26$. وبمتابعة
التبريد تنخفض ذوبانية القصدير في (Pb) (الخط المتقطع): فيترسب قليل من (Sn)
داخل بلورات (Pb).
\end{solution}

\begin{solution}{exo:b2:solid-liquid-diagrams:7}
23.3~\% من NaCl كتلةً: $23.3/76.7 = \qty{0.304}{kg}$ من الملح لكل كيلوغرام من
الماء. وعند \qty{-25}{\degreeCelsius}، تحت الأوتكتيك، لا يمكن أن يوجد أي سائل
أيًّا كانت النسب: فيبقى الملح والجليد متجاورين جسمين صلبين.
\end{solution}

\begin{solution}{exo:b2:solid-liquid-diagrams:8}
يقع $x_{\mathrm B} = 0.5$ بين E$_1$ والمركب: فيتبلور $\mathrm{AB_2}$
أولًا، وينتقل السائل نحو E$_1$، ويتجمد هناك إلى
أوتكتيك A و$\mathrm{AB_2}$؛ وفي النهاية، A + $\mathrm{AB_2}$.
ويقع $x_{\mathrm B} = 0.9$ بين E$_2$ وB: فيتبلور B أولًا، ويبلغ السائل
E$_2$؛ وفي النهاية، $\mathrm{AB_2}$ + B.
\end{solution}

\begin{solution}{exo:b2:solid-liquid-diagrams:9}
عند الحافة الخلفية للمنطقة يتجمد السائل إلى جسم صلب لا يحوي إلا 0.78
من كسره المولي للنحاس: فيبقى النحاس المرفوض في السائل، الذي
يتقدم ويغتني، فيُكنس النحاس نحو طرف القضيب.
ومع نسبة قريبة من 1 لا يزيل كل تمرير إلا جزءًا من الشائبة:
فتلزم تمريرات كثيرة، يبدأ كل منها من القضيب الذي تركه سابقه.
\end{solution}

\begin{solution}{exo:b2:solid-liquid-diagrams:10}
الاستنتاج كما في الفصل. ومن أجل سائل مخفف، $\ln x_A \approx -x_B$، $1/T -
1/T^* \approx -\Delta T/T^{*2}$، ومنه $x_B = \Delta_{\text{انصهار}}H_A\Delta T/(RT^{*2})$؛
ومع $x_B \approx n_B/n_A = bM_A$، $\Delta T = (RT^{*2}M_A/\Delta_{\text{انصهار}}H_A)\,b$.
البنزين: $K_f = 8.314 \times 278.64^2 \times 0.07811/9870 = \qty{5.11}{K.kg/mol}$.
\end{solution}

\begin{solution}{exo:b2:solid-liquid-diagrams:11}
لكل \qty{100}{g}: $19.0/24.31 = \qty{0.782}{mol}$ من Mg، و$81.0/207.2 =
\qty{0.391}{mol}$ من Pb؛ والنسبة $2.00$: $\ce{Mg2Pb}$.
\end{solution}

\begin{solution}{exo:b2:solid-liquid-diagrams:12}
يتناسب الاستواء مع كتلة السائل الأوتكتيكي: فالسبيكة المجهولة تحوي
$120/300 = 0.40$ منها. ففي جهة الرصاص، $w = 18.91 + 0.40 \times (62.13 -
18.91) = 36.2~\%$ من القصدير؛ وفي جهة القصدير، $w = 96.59 - 0.40 \times (96.59 - 62.13) =
82.8~\%$. ويميّز بينهما أول انكسار في \omterm{def:b2:solid-liquid-diagrams:cooling-curve}{منحنى التبريد} (الأعلى للسبيكة ذات 36~\%،
التي يقع خط سيولتها على جهة الرصاص الشديدة الانحدار) أو صورة مجهرية (بلورات
أولية من (Pb) أو من (Sn)).
\end{solution}

\begin{solution}{pb:b2:solid-liquid-diagrams:1}
\textbf{1.} الرصاص \qty{327.5}{\degreeCelsius}، والقصدير \qty{231.9}{\degreeCelsius}.
\textbf{2.} لكل \qty{100}{g}: $62.13/118.71 = \qty{0.5234}{mol}$ من القصدير 
و$37.87/207.2 = \qty{0.1828}{mol}$ من الرصاص؛ $x_{\mathrm{Sn}} = 0.5234/0.7062 =
0.741$.
\textbf{3.} النقاط (0، 327.5)، (100، 231.9)، (62.13، 182.2)، (18.91، 182.2)،
(96.59، 182.2).
\textbf{4.} كما في الشكل: سائل؛ L + (Pb)؛ L + (Sn)؛ (Pb)؛ (Sn)؛ (Pb) + (Sn).
\textbf{5.} \babelsublr{L (62.1~\% Sn) $\longrightarrow$ (Pb) (18.9~\%) + (Sn) (96.6~\%)}.
\textbf{6.} $v = 2 + 1 - 3 = 0$: فدرجة الحرارة والتراكيب الثلاثة
محددة.
\textbf{7.} \omterm{def:b2:solid-liquid-diagrams:solid-solution}{المحلول الصلب} الغني بالرصاص (Pb): فالرصاص يذيب القصدير، فيحوي الجسم الصلب في
توازن مع السائل قصديرًا منذ البداية، بالنسبة المقروءة على
\omterm{def:b2:solid-liquid-diagrams:liquidus}{خط الصلابة}.
\textbf{8.} يتبع \omterm{def:b2:solid-liquid-diagrams:liquidus}{خط السيولة} نحو الأوتكتيك، مغتنيًا بالقصدير،
حتى 62.13~\%.
\textbf{9.} سائل فيه 62.13~\% و(Pb) فيه 18.91~\% من القصدير.
\textbf{10.} (Pb): $(62.13 - 40)/(62.13 - 18.91) = 0.512$؛ والسائل 0.488.
\textbf{11.} (Pb) فيه 18.91~\% و(Sn) فيه 96.59~\%؛ (Sn): $(40 - 18.91)/(96.59 -
18.91) = 0.271$؛ (Pb): 0.729.
\textbf{12.} المكوِّن الأوتكتيكي 48.8~\% (السائل الموجود فوقه مباشرة)؛
وبلورات (Pb) الأولية 51.2~\%. و(Pb) في السؤال 11 مجموع البلورات
الأولية وصفائح (Pb) في الأوتكتيك.
\textbf{13.} انكسار عند \omterm{def:b2:solid-liquid-diagrams:liquidus}{خط السيولة}، ثم انخفاض أبطأ، ثم استواء عند
\qty{182.2}{\degreeCelsius} مدته 0.488 من مدة استواء السبيكة الأوتكتيكية.
\textbf{14.} $\ln x_{\mathrm{Sn}} = -(7030/8.314)(1/T - 1/505.1)$.
\textbf{15.} $\ln x_{\mathrm{Bi}} = -(11\,300/8.314)(1/T - 1/544.6)$.
\textbf{16.} $x_{\mathrm{Sn}}(T) + x_{\mathrm{Bi}}(T) = 1$.
\textbf{17.} عند \qty{110}{\degreeCelsius}: $0.587 + 0.349 = 0.936$؛ وعند
\qty{120}{\degreeCelsius}: $0.621 + 0.382 = 1.003$؛ وعند
\qty{130}{\degreeCelsius}: $0.655 + 0.417 = 1.072$. ويبلغ المجموع 1 قرب
\qty{119.5}{\degreeCelsius}: أي \qty{120}{\degreeCelsius} بدقة درجة واحدة.
\textbf{18.} $x_{\mathrm{Bi}} = 0.38$؛ وكتلةً $0.38 \times 208.98/(0.38 \times
208.98 + 0.62 \times 118.71) = 0.52$.
\textbf{19.} يعطي النموذج \qty{120}{\degreeCelsius} و52~\% من البزموت، والمخطط
المقيَّم \qty{138.8}{\degreeCelsius} و57~\%: أي أدنى بنحو
\qty{19}{\degreeCelsius}.
\textbf{20.} القصدير الذي يتبلور ليس نقيًا بل محلول يحوي
البزموت: فجهده الكيميائي أدنى من جهد القصدير النقي، والجسم الصلب
أكثر استقرارًا، فيتبلور من السائل عند درجة حرارة أعلى من أجل
تركيب معيّن للسائل. فيرتفع فرع القصدير من \omterm{def:b2:solid-liquid-diagrams:liquidus}{خط السيولة}، ويلاقي
فرع البزموت أعلى (والسائل غير المثالي يضيف تصحيحه الخاص).
\textbf{21.} ينصهر ويتجمد عند درجة حرارة واحدة، هي الأدنى في النظام:
فلا مجال عجينيًا قد تتشقق خلاله وصلة تحركت.
\textbf{22.} الرصاص سام، والنفايات الإلكترونية تنتهي في البيئة.
\textbf{23.} يمكن لحام المكوّنات الحساسة للحرارة عند درجة حرارة أدنى؛
لكن الوصلة تلين عند درجة حرارة أدنى في الاستعمال، والبزموت يجعلها
هشة.
\textbf{24.} \qty{570}{g} من البزموت و\qty{430}{g} من القصدير.
\textbf{25.} يتنبأ نموذج السائل المثالي بأوتكتيك عند $\boldsymbol{T_E
\approx \qty{120}{\degreeCelsius}}$، مقابل \qty{138.8}{\degreeCelsius}
المقيَّمة.
\end{solution}
